\chapter{Asymmetric Channels and Metric Dependence}
\label{app:asymmetric}

\section{General weighted cancellation theorem}

Let $a,b\in\C\setminus\{0\}$ and define
\[
 \Acomp_{a,b}(\sigma,\phi)
 =a\sqrt\sigma+b e^{i\phi}\sqrt{1-\sigma}.
\]

\begin{theorem}[Weighted cancellation]
\label{thm:weighted-cancellation}
For $0<\sigma<1$, $\Acomp_{a,b}(\sigma,\phi)=0$ if and only if
\begin{equation}
 \sigma=\frac{|b|^2}{|a|^2+|b|^2}
 \label{eq:weighted-sigma-app}
\end{equation}
and
\begin{equation}
 e^{i\phi}=-\frac{a}{b}\frac{|b|}{|a|}.
 \label{eq:weighted-phase-app}
\end{equation}
The phase on the right has unit modulus.
\end{theorem}
\begin{proof}
The zero equation implies
\[
 |a|\sqrt\sigma=|b|\sqrt{1-\sigma}.
\]
Squaring and solving gives \cref{eq:weighted-sigma-app}. Dividing the original equation by $b\sqrt{1-\sigma}$ then gives \cref{eq:weighted-phase-app}. The converse is direct.
\end{proof}

\begin{corollary}
The cancellation point is $\sigma=1/2$ if and only if $|a|=|b|$.
\end{corollary}

\section{Metric interpretation}

A detector is a covector $d=(\bar a,\bar b)$ acting on a state. Equal gain means the detector has equal modulus on the two basis vectors. This condition depends on the chosen inner product and basis. If the channel basis is changed by a non-unitary transformation, equal coefficients do not represent equal physical or geometric gain.

Therefore a canonical zeta-state construction must provide:
\begin{enumerate}
  \item a Hilbert-space inner product;
  \item an orthogonal complement-related decomposition;
  \item a swap isometry $X$;
  \item a detector invariant under the swap.
\end{enumerate}
If $X$ is unitary and $X\ket0=\ket1$, then swap invariance of the detector ray implies equal gain up to phase.

\section{Perturbation of the selected line}

Write $|a|=1+\epsilon$ and $|b|=1-\epsilon$ for small $\epsilon$. Then
\begin{align*}
 \sigma_*(\epsilon)
 &=\frac{(1-\epsilon)^2}{(1+\epsilon)^2+(1-\epsilon)^2}\\
 &=\frac12-\epsilon+O(\epsilon^3).
\end{align*}
A first-order gain asymmetry causes a first-order horizontal shift. The half-axis is therefore highly diagnostic of exact channel symmetry.

If a proposed numerical map produces weights close to, but not exactly, $1/2$ at known zeros, it may be detecting a small metric asymmetry rather than numerical noise. The source of that asymmetry must be identified.

\section{Monotone reparameterizations}

The direct ansatz uses channel probability $p(\sigma)=\sigma$. A more general complement-compatible map satisfies
\begin{equation}
 p(1-\sigma)=1-p(\sigma).
 \label{eq:complement-p}
\end{equation}
If $p$ is strictly monotone, then $p(1/2)=1/2$ and the balanced line remains unchanged. Thus the location theorem does not require the identity map specifically; it requires a complement-symmetric injective weight map.

\begin{proposition}[Generalized weight theorem]
Let $p:(0,1)\to(0,1)$ be injective and satisfy \cref{eq:complement-p}. If an equal-gain readout with weights $p(\sigma)$ and $1-p(\sigma)$ vanishes, then $\sigma=1/2$.
\end{proposition}
\begin{proof}
Cancellation gives $p(\sigma)=1/2$. The complement property gives $p(1/2)=1/2$. Injectivity implies $\sigma=1/2$.
\end{proof}

This generalization may be useful if a canonical metric produces a nonlinear probability coordinate. The exact identity $p(\sigma)=\sigma$ is conceptually clean but not logically necessary for the half-axis conclusion.

\section{Indefinite metrics}

If the channel space has an indefinite inner product, equal ``norms'' can lose positivity and cancellation conditions can change dramatically. Since the information interpretation requires probabilities and the detector proof uses positivity, the initial ansatz restricts to positive-definite Hilbert metrics. Any indefinite or Krein-space extension must replace the entropy and positivity arguments with new ones.

\status{The weighted-channel theorem and monotone reparameterization result are exact. They expose metric symmetry as a required, testable assumption.}
